\SourceLeaf{251}%
\Needspace{7\baselineskip}
\begin{center}
\textsc{Gauss} to \textsc{Gerling.} Göttingen, 17 April 1844.
\end{center}

The way you prove equality of volume for solids that are merely symmetric and lack congruence gave me much pleasure. One could bring out the essential point by saying

1) that one calls attention to the pyramids whose symmetric counterparts are also congruent to them, namely those in which two faces are normal to a third and form equal edges where they intersect it, and

2) that one shows how every pyramid can be divided into 12 pyramids of the kind just described.

Whether the method of proof is new I cannot, however, guarantee with certainty. At least consult \textsc{Legendre}'s Geometry, in which he attempts to prove a number of things more rigorously or more simply than has usually been done. I do not have the book at hand just now and do not own a copy myself.

Now that this proposition no longer falls under it, I must confine my \textls{regret} to the other propositions of solid geometry that still depend on the method of exhaustion, such as Euclid XII. 5. Perhaps there is room for improvement here too. At present I do not have time to give the subject further thought. \mbox{.\ .\ .\ .\ .\ .}

\Needspace{7\baselineskip}
\begin{center}
[\textsc{Gerling} \emph{to} \textsc{Gauss}. \emph{Marburg, 7 July 1844.}]
\end{center}

\{\mbox{.\ .\ .\ .\ .\ .\ .} First, as to symmetry, where I gave only the “fact” in my next-to-last letter, the \textls{reason} seems to me to lie in our intuition of the straight line and plane. By virtue of it, in the plane we can conceive only \textls{one} perpendicular to a straight line at a point, but with this perpendicular extended to infinity on \textls{both sides}, and likewise\space%
